% !TEX root = ../main.tex

\section{Úvod}

Regulácia teploty kvapalín je dôležitá v mnohých technických oblastiach, napríklad 
pri chladení elektroniky, laboratórnych meraniach, technologických procesoch alebo 
testovaní tepelných sústav. Cieľom takýchto systémov je udržiavať teplotu média na 
požadovanej hodnote a podľa nameraných údajov riadiť akčné členy.

V rámci tímového projektu bol navrhnutý inteligentný IoT systém na reguláciu teploty 
kvapaliny v uzavretom hydraulickom okruhu. Ako hlavný chladiaci prvok je použitý 
Peltierov článok, ktorého studená strana odoberá teplo z kvapaliny cez výmenník tepla. 
Teplá strana článku musí byť aktívne chladená chladičom a ventilátorom.

Systém je doplnený o mikrokontrolér ESP32, ktorý zabezpečuje meranie teploty a ovládanie 
výkonových prvkov. Namerané údaje sú odosielané na Flask server, kde sa spracúvajú, 
ukladajú a zobrazujú vo webovom rozhraní. Súčasťou riešenia je aj integrácia s platformou 
ThingsBoard, ktorá umožňuje vzdialený monitoring a riadenie systému.

\subsection{Cieľ projektu}

Cieľom projektu je navrhnúť, zostaviť a oživiť laboratórny model inteligentného systému na 
reguláciu teploty kvapaliny pomocou Peltierovho článku. Systém má umožňovať meranie aktuálnej 
teploty, riadenie chladiaceho výkonu, ovládanie prietoku kvapaliny a vzdialené monitorovanie 
nameraných údajov.

Výsledkom projektu má byť funkčný laboratórny prototyp, ktorý prepája hardvérovú časť, 
riadiaci softvér, lokálne webové rozhranie a cloudovú IoT platformu.

\subsection{Zadanie projektu}

Tímové zadanie bolo zamerané na realizáciu laboratórneho modelu na reguláciu teploty 
chladiacej kvapaliny v uzavretom hydraulickom okruhu. Z hardvérového hľadiska bolo potrebné 
vytvoriť vodný okruh s nádržkou, pumpou, hadičkami a výmenníkom tepla pripojeným 
k studenej strane Peltierovho článku. Súčasťou systému mal byť aj snímač teploty, 
výkonové rozhranie pre riadenie akčných členov a vyhrievacie teleso na simuláciu tepelnej záťaže.

Softvérová časť zadania zahŕňala čítanie teploty, generovanie PWM signálov pre pumpu a 
Peltierov článok, komunikáciu medzi mikrokontrolérom a serverom, vytvorenie webového rozhrania, 
ukladanie dát a integráciu s platformou ThingsBoard.

Keďže projekt realizoval sedemčlenný tím, rozsah zadania zahŕňal aj identifikáciu systému, 
realizáciu troch regulačných módov, integráciu do ThingsBoard dashboardu a spracovanie 
dokumentácie. Regulačné módy boli definované nasledovne:

\begin{itemize}
    \item \textbf{Mód 1} -- regulácia pomocou Peltierovho článku pri konštantnom výkone pumpy,
    \item \textbf{Mód 2} -- regulácia pomocou prietoku pumpy pri konštantnom výkone Peltierovho článku,
    \item \textbf{Mód 3} -- hybridná alebo kaskádová regulácia využívajúca Peltierov článok aj pumpu.
\end{itemize}

